\documentclass[preprint,12pt]{elsarticle}

\usepackage{amssymb}
\usepackage{amsmath}
\usepackage{hyperref}
\usepackage{lineno}
\usepackage{xcolor}
\usepackage{array}
\usepackage{multirow}
\usepackage{pifont}
\usepackage{adjustbox}
\usepackage{booktabs}
\usepackage{rotating}
\usepackage{colortbl}
\usepackage{tabularx}
\usepackage{makecell}
\usepackage{threeparttable}
\usepackage{soul}

\newcommand{\cmark}{\ding{51}}
\newcommand{\xmark}{\ding{55}}
\newcommand{\pmark}{\textcolor{orange}{$\sim$}}

\newcommand{\rev}[1]{\textcolor{blue}{#1}}
\newcommand{\del}[1]{\textcolor{red}{\sout{#1}}}

\journal{Machine Vision and Applications}

\begin{document}

\begin{frontmatter}

%%---------------------------------------------------------------------------
%%  TITLE
%%---------------------------------------------------------------------------
\title{A Transfer Learning Framework for Pneumonia Detection from Chest X-Ray Images Using Custom CNN and MobileNetV2}

%%---------------------------------------------------------------------------
%%  AUTHOR
%%---------------------------------------------------------------------------
\author[label1]{Gajula Tharun Kumar}

\affiliation[label1]{%
  organization={Department of Data Science, University of Europe for Applied Sciences},
  addressline={Think Campus, Konrad-Zuse-Ring 11},
  city={Potsdam},
  postcode={14469},
  state={Brandenburg},
  country={Germany}}

%%---------------------------------------------------------------------------
%%  ABSTRACT
%%---------------------------------------------------------------------------
\begin{abstract}
Pneumonia is a life-threatening respiratory infection responsible for
millions of deaths globally each year, disproportionately affecting
children under five and elderly patients in low-resource settings.
Timely and accurate diagnosis from chest radiographs is essential yet
remains constrained by radiologist availability, inter-observer
variability, and high workload in clinical environments.

Despite substantial progress in deep-learning-based medical image
classification, most existing studies rely on a single model, neglect
class imbalance, and provide no visual explanation of model decisions,
limiting clinical trustworthiness.
This study addresses the gap by developing a comparative convolutional
neural network framework for binary pneumonia classification from chest
X-ray images.

The publicly available Chest X-Ray Images (Pneumonia) dataset from
Kaggle, containing 5{,}863 images across two classes (NORMAL and
PNEUMONIA), was used.
A custom three-block CNN was trained as a baseline and compared with a
MobileNetV2 transfer-learning model pre-trained on ImageNet.
Both models were evaluated on an independent held-out test set using
accuracy, precision, recall, and F1-score.

MobileNetV2 achieved a test accuracy of 88.46\% and a test loss of
0.2419, outperforming the custom CNN (71.47\% accuracy, loss 0.5932).
Transfer learning reduced test loss by 59\% and improved generalization
considerably, particularly for pneumonia recall (0.96 versus 0.55).

This work delivers a reproducible and computationally efficient CNN
pipeline suitable for pneumonia screening support in resource-constrained
clinical settings.
Grad-CAM and SHAP explainability analysis are identified as important
directions for future work.
\end{abstract}

%%---------------------------------------------------------------------------
%%  HIGHLIGHTS
%%---------------------------------------------------------------------------
\begin{highlights}
\item Pneumonia detection from chest X-rays poses critical diagnostic
      challenges in low-resource clinical settings.
\item A custom CNN and MobileNetV2 are trained and compared on 5{,}863
      chest X-ray images across two classes.
\item MobileNetV2 achieves 88.46\% accuracy, outperforming the custom
      CNN (71.47\%) by 17 percentage points.
\item MobileNetV2 achieves 0.96 recall on PNEUMONIA, a clinically
      critical metric for screening support.
\item The proposed framework provides an efficient, reproducible pipeline
      deployable in resource-constrained environments.
\end{highlights}

%%---------------------------------------------------------------------------
%%  KEYWORDS
%%---------------------------------------------------------------------------
\begin{keyword}
Pneumonia detection \sep
Chest X-ray classification \sep
Convolutional neural network \sep
Transfer learning \sep
MobileNetV2 \sep
Deep Learning
\end{keyword}

\end{frontmatter}

\linenumbers

%%===========================================================================
\section{Introduction}
%%===========================================================================
The Introduction establishes the clinical problem, motivates the use of
CNN-based image analysis, identifies the research gap, and states the
exact contributions of this study.

%%---------------------------------------------------------------------------
\subsection{Application Background and Practical Importance}
%%---------------------------------------------------------------------------
Pneumonia is an acute lower respiratory infection caused by bacteria,
viruses, or fungi that inflame the air sacs in one or both lungs,
causing them to fill with fluid or pus.
It is among the leading causes of death worldwide, responsible for
approximately 2.5 million fatalities annually, with the burden
concentrated in sub-Saharan Africa and South Asia~\cite{reshan2023}.
Children under the age of five account for over 700{,}000 pneumonia
deaths per year, making it the single largest infectious cause of child
mortality globally~\cite{liang2020}.
Early and accurate diagnosis is essential for timely antibiotic or
antiviral treatment, yet access to trained radiologists is severely
limited in many high-burden regions~\cite{rajpurkar2017chexnet}.
The gold standard for pneumonia diagnosis is expert interpretation of
chest radiographs, which requires considerable clinical skill and
introduces significant inter-observer variability~\cite{hashmi2020}.
In high-volume emergency settings, radiologists may be required to
interpret hundreds of images per shift, increasing the risk of
diagnostic error~\cite{chouhan2020}.
Automated image analysis systems capable of reliably flagging pneumonia
on chest X-rays could reduce workload, accelerate triage, and extend
diagnostic reach to underserved populations~\cite{rajpurkar2017chexnet}.
Such systems are particularly valuable in resource-constrained hospitals
where 24-hour radiologist coverage is unavailable.

%%---------------------------------------------------------------------------
\subsection{Machine Learning and Computer Vision Context}
%%---------------------------------------------------------------------------
Convolutional neural networks (CNNs) have transformed medical image
analysis by learning hierarchical visual representations directly from
raw pixel data without hand-crafted feature engineering~\cite{krizhevsky2012}.
Successive convolutional layers detect progressively abstract patterns,
from edges and textures in early layers to organ-level structures in
deeper layers, making them well suited to radiographic analysis.
Unlike traditional machine-learning pipelines that require domain
experts to define features manually, CNNs adapt their internal
representations to the specific statistics of the training dataset.
Local receptive fields allow CNNs to capture spatially localized
pathological changes such as lung opacities, consolidations, and
infiltrates that are hallmarks of pneumonia on chest X-rays.
The spatial parameter sharing of convolutional layers also reduces the
number of learnable parameters compared with fully connected networks,
partially mitigating overfitting on small medical datasets.
However, training a deep CNN from random initialization typically
requires tens of thousands of labeled examples, which are difficult and
expensive to obtain in the medical domain due to the need for expert
annotation and patient privacy constraints~\cite{huang2017densenet}.
These challenges make transfer learning particularly attractive for
medical imaging applications.

%%---------------------------------------------------------------------------
\subsection{Transfer Learning for Pneumonia Detection}
%%---------------------------------------------------------------------------
Transfer learning exploits visual representations learned on large
general-purpose datasets, most commonly ImageNet with over 1.2 million
labeled images, and adapts them to a target domain with limited
labeled data~\cite{liang2020}.
Pre-trained models encode low-level features such as edge detectors and
color gradients that transfer well across visual domains, including
natural images and radiographs.
By freezing the backbone weights and training only a newly attached
classification head, it is possible to achieve strong performance even
when fewer than a few thousand domain-specific labeled images are
available.
Fine-tuning subsequently unfreezes deeper backbone layers and allows
the model to specialize to domain-specific features at a lower learning
rate.
MobileNetV2 is particularly attractive for medical applications because
its inverted residual blocks and linear bottlenecks achieve competitive
accuracy with a fraction of the parameters of heavier architectures
such as VGG16 or InceptionV3, enabling deployment on
resource-constrained hardware~\cite{mobilenetv2}.
The depthwise separable convolutions in MobileNetV2 reduce computation
by factorizing standard convolutions, making inference fast enough for
real-time clinical screening tools.
Selecting the appropriate architecture and fine-tuning strategy
requires systematic comparison, which existing studies in this domain
have not consistently provided.

%%---------------------------------------------------------------------------
\subsection{Deep Learning and Clinical Trustworthiness}
%%---------------------------------------------------------------------------
Accuracy alone is insufficient to establish trust in a clinical
decision-support system; clinicians and regulators require
interpretable evidence that a model reasons correctly~\cite{reshan2023}.
A model may achieve high test accuracy by exploiting spurious
correlations in the training data, such as scanner artifacts, patient
positioning labels, or demographic confounders, rather than true
pathological features~\cite{chouhan2020}.
Gradient-weighted Class Activation Mapping (Grad-CAM) produces
class-specific heatmaps by computing the gradient of the predicted
class score with respect to the final convolutional feature maps,
allowing visualization of which spatial regions drove the prediction.
SHAP (SHapley Additive exPlanations) assigns pixel-level or
superpixel-level contribution scores based on cooperative game theory,
indicating both the magnitude and direction of each region's influence
on the output.
In a medical setting, false negatives carry severe clinical
consequences; a model that misses pneumonia in a critically ill patient
may delay life-saving treatment.
Explainability methods therefore serve not only as debugging tools but
as a fundamental component of the clinical validation pipeline.

%%---------------------------------------------------------------------------
\subsection{Research Gap}
%%---------------------------------------------------------------------------
Despite the large body of work on deep learning for chest X-ray
analysis, several important methodological gaps remain.
Most published studies evaluate only a single model or compare models
without controlling for dataset splitting, augmentation strategy, or
class-imbalance treatment, making cross-study comparisons
unreliable~\cite{rajpurkar2017chexnet,hashmi2020}.
The severe class imbalance present in the standard chest X-ray
pneumonia dataset (approximately 3.9:1 pneumonia-to-normal ratio in
training) is rarely addressed explicitly, risking models that optimize
for majority-class accuracy.
Deep Learning methods are applied in relatively few studies, and where
they are used, explanation quality is rarely assessed
systematically~\cite{reshan2023}.
Computational efficiency is seldom reported despite its importance for
deployment in low-resource hospitals.

%%---------------------------------------------------------------------------
\subsection{Problem Statement}
%%---------------------------------------------------------------------------
This study addresses the problem of automatically classifying chest
X-ray images into NORMAL and PNEUMONIA categories using a comparative
CNN-based framework.
The study investigates whether transfer learning with MobileNetV2 can
improve predictive performance over a custom CNN while maintaining
computational efficiency, and identifies explainability analysis as a
direction for future work.

%%---------------------------------------------------------------------------
\subsection{Aim and Objectives}
%%---------------------------------------------------------------------------
The aim of this study is to develop and evaluate a CNN-based framework
for the reliable binary classification of chest X-ray images into
NORMAL and PNEUMONIA classes.

The specific objectives are:
\begin{itemize}
  \item to analyze and preprocess the Chest X-Ray Images (Pneumonia)
        dataset, including class distribution, image quality, and
        augmentation design;
  \item to develop a custom three-block CNN as a reproducible baseline;
  \item to implement MobileNetV2 with frozen-backbone transfer learning
        and a custom classification head;
  \item to compare both models using accuracy, precision, recall,
        F1-score, loss, training time, and parameter count;
  \item to evaluate training convergence, stability, and signs of
        overfitting on validation curves;
  \item to assess limitations, deployment suitability, and practical
        clinical value of the best-performing model.
\end{itemize}

%%---------------------------------------------------------------------------
\subsection{Research Questions}
%%---------------------------------------------------------------------------
\begin{itemize}
  \item \textbf{RQ1:} How does preprocessing and augmentation affect
        the quality and class balance of the chest X-ray dataset, and
        what are the characteristics of the training, validation, and
        test splits?
  \item \textbf{RQ2:} Which model, the custom CNN or MobileNetV2,
        achieves superior overall predictive performance on the held-out
        test set, and by how much?
  \item \textbf{RQ3:} How reliably does each model classify individual
        classes, and what error patterns appear in the confusion
        matrices and class-level metrics?
  \item \textbf{RQ4:} How do the two models compare in terms of
        computational efficiency, including parameter count, training
        time, and model size, and which is more suitable for deployment
        in resource-constrained clinical environments?
\end{itemize}

%%---------------------------------------------------------------------------
\subsection{Contributions and Novelty}
%%---------------------------------------------------------------------------
This study contributes the following:
\begin{itemize}
  \item A fully reproducible comparative experimental pipeline
        evaluating a custom CNN and MobileNetV2 on the same dataset
        split, preprocessing strategy, and evaluation protocol.
  \item Systematic class-level reliability analysis including per-class
        precision, recall, F1-score, and confusion matrix decomposition,
        enabling identification of clinically high-risk error types.
  \item An efficiency-performance trade-off analysis including parameter
        count, model file size, and test loss, to support deployment
        decision-making in resource-limited settings.
  \item A critically assessed discussion of limitations, dataset bias,
        and practical deployment considerations relevant to clinical
        radiograph screening.
\end{itemize}

%%---------------------------------------------------------------------------
\subsection{Report Organization}
%%---------------------------------------------------------------------------
The remainder of this report is structured as follows.
Section~2 reviews the relevant literature on deep learning for chest
X-ray analysis, transfer learning architectures, preprocessing methods,
and model evaluation.
Section~3 describes the complete methodology including dataset
description, preprocessing, model architectures, training procedure,
and evaluation framework.
Section~4 presents the experimental results organized by research
question.
Section~5 interprets the results, discusses limitations, compares
findings with the literature, and addresses practical and ethical
implications.
Section~6 concludes the study.

%%===========================================================================
\section{Literature Review}
%%===========================================================================

%%---------------------------------------------------------------------------
\subsection{Image-Based Machine Learning in Medical Imaging}
%%---------------------------------------------------------------------------
Machine learning has progressively transformed medical image analysis
over the past two decades, evolving from hand-crafted feature pipelines
to end-to-end learned representations.
Early approaches relied on domain-expert-designed descriptors such as
histogram of oriented gradients, local binary patterns, and Gabor
filters combined with support vector machines or random forests, which
required substantial engineering effort and did not generalize well
across imaging modalities~\cite{krizhevsky2012}.
The introduction of deep convolutional neural networks, popularized by
AlexNet's performance on ImageNet in 2012, demonstrated that
hierarchical feature learning from raw pixel data could surpass
hand-crafted approaches on large-scale visual recognition
tasks~\cite{krizhevsky2012}.
This insight rapidly transferred to the medical domain: CNNs were
applied to dermoscopy, histopathology, ophthalmology, and radiology,
consistently outperforming classical methods when sufficient labeled
data were available.
Chest radiograph analysis has been a particularly active area, driven
by the clinical prevalence of chest X-rays and the availability of
large annotated datasets~\cite{rajpurkar2017chexnet}.
Rajpurkar et al.\ demonstrated that a 121-layer DenseNet could
detect pneumonia on chest X-rays at a level exceeding the average
performance of radiologists on several pathology classes~\cite{rajpurkar2017chexnet}.
Subsequent work confirmed that model performance scales with dataset
size and architecture depth, but also revealed that small or
single-institution datasets introduce overfitting and limit
generalizability~\cite{hashmi2020}.
The reliance on large labeled corpora remains a fundamental constraint
in medical imaging, motivating the adoption of transfer learning,
semi-supervised learning, and data augmentation as practical
mitigations~\cite{reshan2023}.
Recent surveys confirm that transfer learning from ImageNet
consistently outperforms training from scratch on medical imaging
benchmarks of fewer than 10{,}000 labeled samples~\cite{liang2020}.

%%---------------------------------------------------------------------------
\subsection{CNN-Based Methods for Chest X-Ray Classification}
%%---------------------------------------------------------------------------
A substantial body of literature applies CNN-based classification to
the Kermany chest X-ray dataset, and several consistent patterns
emerge across these studies.
Transfer learning with architectures pre-trained on ImageNet is the
dominant approach: VGG16, ResNet50, InceptionV3, DenseNet121, and
MobileNet variants are the most frequently used backbones, and
virtually all reported results from frozen or fine-tuned pre-trained
models outperform custom CNNs trained from
scratch~\cite{chouhan2020,liang2020,reshan2023}.
Reported test accuracies on the binary NORMAL/PNEUMONIA task cluster
between 85\% and 96\% across these studies, but direct comparison is
complicated by inconsistent dataset splitting: some studies use the
original Kaggle split while others reshuffle the data or apply
cross-validation, making reported figures non-comparable~\cite{hashmi2020}.
Class imbalance, present at approximately 2.9:1 in the original
training split, is addressed inconsistently: some studies apply
class-weighted loss or oversampling while others ignore it entirely,
a methodological inconsistency that likely inflates reported accuracy
figures by rewarding majority-class prediction~\cite{reshan2023}.
Augmentation strategies vary widely, with most studies applying
rotation, zoom, and horizontal flip but rarely justifying specific
parameter choices with respect to chest radiograph
anatomy~\cite{chouhan2020}.
Evaluation is frequently limited to overall accuracy, with fewer
studies reporting per-class precision, recall, and F1-score; this
is a significant limitation in a clinical task where false negative
rate for the minority pathology class is the most safety-critical
metric~\cite{hashmi2020}.
Across the reviewed studies, models achieving high overall accuracy
sometimes exhibit substantially degraded recall on the NORMAL class
or vice versa, confirming that single-metric evaluation is
insufficient for clinical assessment.

%%---------------------------------------------------------------------------
\subsection{Transfer Learning Architectures}
%%---------------------------------------------------------------------------
MobileNetV2, introduced by Sandler et al., is built around inverted
residual blocks in which the input and output of each block are
narrow low-dimensional bottleneck layers while the internal expansion
uses a wider high-dimensional representation~\cite{mobilenetv2}.
Within each block, depthwise separable convolutions factorize the
standard convolution into a depthwise spatial filter followed by a
pointwise $1 \times 1$ convolution, reducing the number of multiply-add
operations by a factor of approximately 8 to 9 compared with a
standard convolution of equal depth~\cite{mobilenetv2}.
Linear bottleneck activations are used at the residual connection
endpoints to preserve information in the low-dimensional manifold,
avoiding the information loss caused by applying ReLU to low-rank
representations.
Trained on ImageNet with 1.2 million images and 1000 classes,
MobileNetV2 achieves 72\% top-1 accuracy with only 3.4 million
parameters, making it competitive with substantially heavier models.
When the classification head is removed and the backbone is used as a
feature extractor, the final convolutional feature map encodes rich
multi-scale visual representations that generalize well to domain-
shifted inputs including medical images.

ResNet50~\cite{he2016resnet} introduced residual skip connections that
allow gradients to bypass convolutional blocks, enabling stable
training of networks with 50 or more layers.
Its 25 million parameters provide a larger representational capacity
than MobileNetV2, which yields marginal accuracy gains on large
datasets but increases inference cost and memory consumption.
DenseNet121~\cite{huang2017densenet} connects each layer to all
subsequent layers within a dense block, encouraging feature reuse and
reducing the need for redundant parameter learning; it has been widely
adopted for chest X-ray analysis following its use in CheXNet~\cite{rajpurkar2017chexnet}.
For the present study, MobileNetV2 was selected over ResNet50 and
DenseNet121 because its substantially smaller parameter count and
lower inference cost are better aligned with the goal of deployment
in resource-constrained clinical settings.

%%---------------------------------------------------------------------------
\subsection{Data Preprocessing and Augmentation}
%%---------------------------------------------------------------------------
Consistent preprocessing is essential for both training stability and
fair model comparison.
Resizing all images to a fixed spatial dimension eliminates variability
introduced by different acquisition devices and ensures that the
convolutional receptive field covers a consistent anatomical extent;
$224 \times 224$ pixels is the standard input size for ImageNet-
pretrained architectures including MobileNetV2~\cite{mobilenetv2}.
Pixel normalization to the range $[0, 1]$ or $[-1, 1]$ ensures that
activation magnitudes remain in a regime where gradient-based
optimization is numerically stable, and is standard practice across
CNN-based medical imaging pipelines~\cite{shorten2019}.
Data augmentation artificially expands the effective training set size
by applying stochastic geometric and photometric transformations,
reducing overfitting when the labeled dataset is small.
In the context of chest radiographs, clinically plausible
augmentations include small rotations to simulate patient positioning
variability, zoom to account for differences in field of view and
patient size, and width and height shifts to tolerate variation in
lung centering~\cite{shorten2019}.
Horizontal flip is appropriate for frontal chest radiographs because
the overall lung pattern is approximately symmetric and the left-right
orientation does not reliably distinguish pathology from health.
Vertical flip is contraindicated for chest X-rays because the
superior-inferior anatomical orientation is diagnostically meaningful:
the lung apices, bases, costophrenic angles, and cardiac silhouette
each carry distinct diagnostic information that a vertically flipped
image would misrepresent~\cite{chouhan2020}.
Color jitter and grayscale inversion are generally avoided in
radiograph augmentation because radiographic intensity values encode
tissue density and inverting them removes physically meaningful
contrast relationships.
The appropriate magnitude of augmentation parameters should be
calibrated to the expected range of clinical variation; overly
aggressive augmentation introduces unrealistic image distortions
that impair rather than improve generalization.

%%---------------------------------------------------------------------------
\subsection{Model Evaluation and Reliability}
%%---------------------------------------------------------------------------
The choice of evaluation metric has significant consequences for how
model performance is interpreted, particularly in clinically imbalanced
binary classification tasks.
Overall accuracy, defined as the proportion of correctly classified
samples, is sensitive to class imbalance: a model that predicts the
majority class for every input will achieve high accuracy on a dataset
where one class predominates, yet provide no clinical utility~\cite{hashmi2020}.
Precision and recall decompose predictive performance at the class
level: precision measures the fraction of positive predictions that
are correct, while recall (sensitivity) measures the fraction of true
positive cases that are detected.
In pneumonia screening, recall for the PNEUMONIA class is the
safety-critical metric because a false negative represents a missed
diagnosis; a low-recall model would fail patients who require
immediate treatment~\cite{reshan2023}.
The F1-score is the harmonic mean of precision and recall and provides
a single class-level summary that penalizes extreme imbalances between
the two, making it more informative than accuracy alone on imbalanced
datasets.
Specificity complements sensitivity by quantifying the rate of correct
normal classifications, which is relevant for managing unnecessary
referrals in a screening context.
Area under the receiver operating characteristic curve (AUC)
characterizes model discrimination across all possible classification
thresholds and is threshold-independent, making it useful for comparing
models whose optimal operating threshold may differ.
The confusion matrix provides the complete joint distribution of
predicted and true labels and is the most transparent summary of
classification behavior, enabling identification of which specific
error types (false positives or false negatives) dominate model
failure~\cite{hashmi2020}.
Several reviewed studies report only overall accuracy, which obscures
class-level error patterns and risks overstating clinical utility;
per-class metrics and confusion matrices should be considered the
minimum reporting standard for medical classification models.

%%---------------------------------------------------------------------------
\subsection{Deep Learning in Medical Image Analysis}
%%---------------------------------------------------------------------------
Deep Learning methods for image classification generate visual or
numerical evidence of which input regions or features drove a model's
prediction, addressing the opacity of deep neural networks that
otherwise function as black boxes~\cite{reshan2023}.
In medical imaging, explainability serves both a quality-assurance
function, confirming that the model attends to anatomically relevant
structures, and a regulatory function, providing the evidence required
for clinician trust and institutional adoption.

Gradient-weighted Class Activation Mapping (Grad-CAM) computes a
class-discriminative localization map by taking the global average of
the gradients of the predicted class score with respect to the
activations of a target convolutional layer, then forming a weighted
sum of those feature maps~\cite{chouhan2020}.
The resulting heatmap is upsampled to input resolution and overlaid
on the original image, highlighting spatial regions that contributed
most to the prediction.
Grad-CAM is widely applied to chest X-ray models because it requires
no architectural modification and is computationally inexpensive;
however, its spatial resolution is limited by the stride of the target
convolutional layer, so fine-grained localization of small lesions
may be imprecise~\cite{reshan2023}.
In chest X-ray studies, Grad-CAM maps have been used to confirm that
correctly classified pneumonia cases show activation in the lung
fields, while misclassified cases sometimes show spurious activation
in background regions, suggesting shortcut learning.

SHAP (SHapley Additive exPlanations) assigns a contribution score to
each feature based on its average marginal contribution across all
possible subsets of features, grounded in cooperative game theory.
For image models, DeepSHAP or GradientSHAP propagate SHAP values
through the network layers, producing pixel-level or superpixel-level
attribution maps that indicate both the magnitude and direction of
each region's influence on the output~\cite{reshan2023}.
Unlike Grad-CAM, SHAP values can be negative, identifying regions
that suppressed the predicted class probability, which provides
richer diagnostic information.
The principal limitation of SHAP for deep image models is
computational cost: exact Shapley value computation is exponential
in the number of features, and approximations require careful
selection of background reference samples.

Despite the availability of these methods, their application in chest
X-ray studies remains inconsistent.
Heatmaps that appear anatomically plausible do not guarantee correct
clinical reasoning; a model may highlight lung regions for the wrong
reason, and agreement between Grad-CAM and SHAP does not constitute
clinical validation~\cite{hashmi2020}.
Domain-expert radiologist assessment of explanation outputs is the
gold standard for evaluating explanation quality, yet this step is
absent from the majority of published studies.

%%---------------------------------------------------------------------------
\subsection{Robustness, Generalization, and Deployment}
%%---------------------------------------------------------------------------
Generalization to unseen data distributions is a fundamental challenge
for medical imaging models trained on single-institution datasets.
Performance degradation under domain shift, caused by differences in
scanner manufacturer, acquisition protocol, patient demographics, or
image processing pipeline, has been documented across multiple medical
imaging tasks and is a principal obstacle to real-world
deployment~\cite{hashmi2020}.
Models trained exclusively on the Kaggle chest X-ray dataset, which
originates from a single pediatric hospital in Guangzhou, China, may
not generalize to data from adult populations or different geographic
regions without retraining or domain adaptation.
Robustness to common image degradations such as Gaussian noise,
motion blur, and low contrast has been studied in natural image
classification but is rarely evaluated in chest X-ray
studies~\cite{reshan2023}.
Such degradations are clinically relevant because radiograph quality
varies with equipment calibration, patient cooperation, and technical
factors, and a model deployed in a resource-constrained setting may
encounter images of lower quality than those in the training set.

Deployment feasibility depends critically on model size, inference
latency, and hardware requirements.
Lightweight architectures such as MobileNetV2 were specifically
designed for inference on mobile and embedded processors, and their
depthwise separable convolution blocks reduce multiply-add operations
sufficiently for real-time inference on CPU-only
hardware~\cite{mobilenetv2}.
By contrast, larger models such as VGG16 (138 million parameters) or
ResNet50 (25 million parameters) require GPU acceleration for
acceptable inference speed, limiting their applicability in hospitals
with no dedicated AI hardware.
Model file size also matters for deployment: quantized MobileNetV2
models compressed to 8-bit integer weights can be stored and loaded
on devices with limited flash memory.
Few chest X-ray classification studies report inference latency or
model file size, leaving deployment feasibility as an unaddressed
gap in the majority of the published literature.

%%---------------------------------------------------------------------------
\subsection{Comparative Literature Review Matrix}
%%---------------------------------------------------------------------------

Table~\ref{tab:literature_matrix}Table summarizes representative studies on deep-learning-based chest X-ray
classification.
The comparison reveals that most studies use the same Kaggle
chest X-ray dataset but differ substantially in model choice,
augmentation strategy, and evaluation protocol.
Transfer learning is widely adopted, yet fewer than half of the
reviewed studies apply any class-imbalance treatment.
Explainability methods are present in only a minority of studies.
Robustness testing and external validation are almost universally
absent, and statistical testing between models is rarely reported.
Efficiency analysis is present in some lightweight-model studies but
absent from most classification papers.
The proposed study directly addresses these gaps by combining a custom
CNN baseline, MobileNetV2 transfer learning, class-weight balancing,
efficiency reporting, and class-level error analysis within a single
reproducible pipeline.

\begin{sidewaystable}[p]
\centering
\tiny
\caption{Comparative literature review matrix summarizing datasets,
models, evaluation practices, explainability, robustness, and
validation strategies in closely related chest X-ray classification
studies using deep learning. Symbols: \cmark~feature explicitly
included; \xmark~feature not included; \pmark~partially included or
unclear. TL: Transfer Learning; CB: Class Balancing; MC:
Multiple-Model Comparison; GC: Grad-CAM; SH: SHAP; RT: Robustness
Testing; EV: External Validation; ST: Statistical Testing; EA:
Efficiency Analysis.}
\label{tab:literature_matrix}
\begin{tabular}{p{1.5cm} p{1.7cm} p{1.3cm} c p{1.1cm} p{1.7cm}
                c c c c c c c c c p{1.6cm} p{2.3cm} p{2.3cm}}
\toprule
\textbf{Study} & \textbf{Problem} & \textbf{Dataset} &
\textbf{Classes} & \textbf{Size} & \textbf{Models} &
\textbf{TL} & \textbf{CB} & \textbf{MC} & \textbf{GC} &
\textbf{SH} & \textbf{RT} & \textbf{EV} & \textbf{ST} &
\textbf{EA} & \textbf{Metrics} & \textbf{Main Contribution} &
\textbf{Principal Limitation} \\
\midrule

Chouhan et al.\ \cite{chouhan2020}, 2020 &
  Pneumonia classification &
  Chest X-Ray (Kaggle) & 2 & 5{,}863 &
  VGG16 &
  \cmark & \xmark & \xmark & \cmark & \xmark &
  \xmark & \xmark & \xmark & \xmark &
  Acc, F1 &
  Fine-tuned VGG16 for binary pneumonia &
  No class balancing; no SHAP; single model \\

Liang \& Zheng\ \cite{liang2020}, 2020 &
  Pneumonia detection &
  Chest X-Ray (Kaggle) & 2 & 5{,}863 &
  ResNet50 &
  \cmark & \pmark & \xmark & \cmark & \xmark &
  \xmark & \xmark & \xmark & \pmark &
  Acc, AUC &
  ResNet50 with custom augmentation &
  No SHAP; accuracy-only primary metric \\

Reshan et al.\ \cite{reshan2023}, 2023 &
  Lung disease classification &
  Chest X-Ray (Kaggle) & 2 & 5{,}863 &
  DenseNet121 &
  \cmark & \cmark & \cmark & \cmark & \xmark &
  \xmark & \cmark & \xmark & \xmark &
  AUC per class &
  Multi-label classification with class weights &
  No SHAP; no efficiency analysis \\

\midrule
\textbf{Proposed study} &
  Binary pneumonia classification &
  Chest X-Ray (Kaggle) & 2 & 5{,}863 &
  Custom CNN, MobileNetV2 &
  \cmark & \cmark & \cmark & \xmark & \xmark &
  \xmark & \xmark & \xmark & \cmark &
  Acc, Prec, Rec, F1, Loss &
  Comparative framework with class-level error analysis &
  Single public dataset; no XAI implemented; no external validation \\

\bottomrule
\end{tabular}
\end{sidewaystable}

%%---------------------------------------------------------------------------
\subsection{Synthesis of Literature and Identified Gap}
%%---------------------------------------------------------------------------
The reviewed literature establishes that CNN-based transfer learning
is an effective approach to binary pneumonia classification from chest
radiographs, with MobileNetV2 and related lightweight architectures
offering a favorable balance of accuracy and computational efficiency.
However, four methodological gaps recur consistently across the
reviewed studies and motivate the design of the present work.

First, most studies evaluate models on non-standard or unreported
dataset splits, precluding meaningful cross-study accuracy comparison.
This study addresses the gap by using the original Kaggle train/test
split without reshuffling and reporting all metrics on the fixed
held-out test set.

Second, class-imbalance treatment is inconsistent or absent, leading
to models that optimize majority-class accuracy at the expense of
minority-class recall, which is the clinically critical metric.
This study addresses the gap by applying class-weight compensation
and reporting per-class precision, recall, and F1-score alongside
overall accuracy.

Third, comparative evaluation of a custom CNN baseline against a
transfer-learning model under identical experimental conditions is
uncommon; most studies present only one architecture.
This study addresses the gap by training both models on the same data
split with identical preprocessing, augmentation, optimizer, and
evaluation protocol.

Fourth, efficiency metrics including parameter count, trainable
parameter count, and training duration are rarely reported despite
their importance for deployment in low-resource settings.
This study addresses the gap by recording and reporting these values
for both models.

The absence of Grad-CAM and SHAP explainability analysis in the
current implementation is acknowledged as a limitation and is
identified as the highest-priority direction for future work,
consistent with the consensus in the reviewed literature that
explainability is essential for clinical trust.
The methodology is described in full in Section~3.

%%===========================================================================
\section{Methodology}
%%===========================================================================

%%---------------------------------------------------------------------------
\subsection{Research Design}
%%---------------------------------------------------------------------------
This study follows a quantitative experimental design in which a
custom CNN and a MobileNetV2 transfer-learning model are independently
trained, validated, and evaluated on the same dataset split using
identical preprocessing and evaluation protocols.
The image task is binary classification (NORMAL versus PNEUMONIA) from
frontal chest radiographs.
Computational efficiency is reported alongside predictive performance
to support deployment decision-making.

\begin{figure*}[!ht] \centering \includegraphics[width=\textwidth]{Figures/workflow.png} \caption{Overall workflow of the proposed pneumonia detection system. The workflow consists of dataset collection, preprocessing, data augmentation, model development using a Custom CNN and MobileNetV2, model training, evaluation using classification metrics and confusion matrices, and final performance comparison.} \label{fig:workflow} \end{figure*}

%%---------------------------------------------------------------------------
\subsection{Dataset Description}
%%---------------------------------------------------------------------------

\subsubsection{Dataset Source and Accessibility}
The Chest X-Ray Images (Pneumonia) dataset was originally published
by Kermany et al.~\cite{kermany2018} and is publicly accessible on
Kaggle at
\url{https://www.kaggle.com/datasets/paultimothymooney/chest-xray-pneumonia}.
The dataset was accessed in 2024 under a CC BY 4.0 license and used
without modification to the original class labels.
It was selected because it is the most widely used public benchmark
for binary pneumonia classification from chest radiographs, enabling
direct comparison with existing literature.

\subsubsection{Dataset Composition}
The dataset contains 5{,}863 chest X-ray images in JPEG format across
two classes: NORMAL and PNEUMONIA.
The images are organized into pre-defined training, validation, and
test folders, with 5{,}216 images in training, 16 images in validation,
and 624 images in test.
Training images consist of 1{,}341 NORMAL and 3{,}875 PNEUMONIA
samples, representing a class imbalance ratio of approximately 2.9:1.
Test images consist of 234 NORMAL and 390 PNEUMONIA samples.
Image resolutions vary across the dataset; all images are resized to
$224 \times 224$ pixels for model input.
Images are grayscale radiographs but are loaded as three-channel RGB
to match the expected input format of MobileNetV2.

\subsubsection{Class Definitions}
The NORMAL class contains chest radiographs from individuals with no
evidence of pulmonary infection, showing clear lung fields without
consolidation, infiltrates, or ground-glass opacities.
The PNEUMONIA class contains radiographs exhibiting pathological
findings consistent with bacterial or viral pneumonia, including
patchy consolidation, air bronchograms, and diffuse bilateral
infiltrates; the class includes both bacterial and viral subtypes
without a separate sub-label.

\subsubsection{Dataset Quality Assessment}
The training split exhibits a significant class imbalance, with
PNEUMONIA images outnumbering NORMAL images by a ratio of approximately
2.9:1, which could bias a model toward predicting the majority class.
The validation split contains only 16 images, which is insufficient for
reliable epoch-by-epoch model selection; this limitation is inherited
from the original dataset structure.
No corrupt or unreadable image files were identified during initial
loading.
Some images exhibit variation in patient positioning, lung field
cropping, and radiograph exposure, which are representative of
real-world clinical variation.

\subsubsection{Dataset Visualization}
\begin{figure}[!ht] \centering \includegraphics[width=0.8\linewidth]{Figures/dataset_samples.png} \caption{Sample chest X-ray images from the NORMAL and PNEUMONIA classes.} \label{fig:dataset_samples} \end{figure}

%%---------------------------------------------------------------------------
\subsection{Data Preparation}
%%---------------------------------------------------------------------------

\subsubsection{Data Cleaning}
All images were verified to be readable using TensorFlow's
\texttt{ImageDataGenerator}; no corrupted or unreadable files were
identified and no images were excluded.
Class labels were confirmed to match the folder structure provided by
the original dataset.

\subsubsection{Data Splitting}
The original pre-defined training, validation, and test folder
structure was retained.
Training images (5{,}216) were used for model fitting with
augmentation.
The 16-image validation set was used for epoch-level monitoring.
The 624-image test set was treated as the independent held-out
evaluation set and was never used during training or hyperparameter
selection.
No augmented versions of test images were created, preventing data
leakage.

\subsubsection{Image Preprocessing}
All images were resized to $224 \times 224$ pixels to match the input
requirements of MobileNetV2 and to provide a consistent input size for
the custom CNN.
Pixel values were normalized to the range $[0, 1]$ by dividing by 255,
following standard practice for CNN training.
Images were loaded in RGB color mode (three channels) using
\texttt{class\_mode='binary'} to produce binary class labels.
No additional contrast enhancement or histogram equalization was
applied, in order to preserve the natural radiographic intensity
distribution.

\subsubsection{Data Augmentation}
Augmentation was applied exclusively to the training set using
TensorFlow's \texttt{ImageDataGenerator} with the following
transformations and parameter ranges~\cite{shorten2019}:
\begin{itemize}
  \item Rotation: up to $\pm 15^{\circ}$, to simulate variability in
        patient positioning;
  \item Zoom: up to 20\%, to handle variation in magnification;
  \item Width shift: up to 10\% of image width, to tolerate lateral
        positioning variation;
  \item Height shift: up to 10\% of image height, to tolerate
        superior-inferior shift;
  \item Horizontal flip: enabled, as left-right orientation does not
        carry diagnostic meaning in frontal chest radiographs.
\end{itemize}
Vertical flip was not applied because superior-inferior lung anatomy
carries meaningful diagnostic information.
Augmentation was performed on-the-fly during training, so no
augmented images were stored on disk.

\subsubsection{Class-Imbalance Handling}
To address the 2.9:1 training-set imbalance, class weights were
computed as inversely proportional to class frequency and passed to
the \texttt{fit()} function, penalizing misclassification of the
minority NORMAL class more heavily.

%%---------------------------------------------------------------------------
\subsection{Model Development}
%%---------------------------------------------------------------------------

\subsubsection{Custom CNN Baseline}
The custom CNN was designed as a three-block sequential architecture
with increasing filter depth.
The input shape is $(224, 224, 3)$.
Block 1 consists of a \texttt{Conv2D} layer with 32 filters of size
$3 \times 3$ and ReLU activation, followed by \texttt{MaxPooling2D}
with pool size $2 \times 2$.
Block 2 consists of a \texttt{Conv2D} layer with 64 filters and ReLU
activation, followed by \texttt{MaxPooling2D}.
Block 3 consists of a \texttt{Conv2D} layer with 128 filters and ReLU
activation, followed by \texttt{MaxPooling2D} and
\texttt{BatchNormalization}.
The feature maps are then flattened and passed to a fully connected
\texttt{Dense} layer with 128 units and ReLU activation.
A \texttt{Dropout} layer with rate 0.5 regularizes the model.
The output layer is a single \texttt{Dense} unit with sigmoid
activation for binary classification.
The model is compiled with binary cross-entropy loss and the Adam
optimizer at a learning rate of $1 \times 10^{-4}$.
The total parameter count is 11{,}169{,}601 (11{,}169{,}345 trainable).

\begin{figure*}[!ht] \centering \includegraphics[width=\textwidth]{Figures/cnn_architecture.png} \caption{Architecture of the proposed Custom CNN model used for pneumonia detection.} \label{fig:cnn_architecture} \end{figure*}

\subsubsection{MobileNetV2 Transfer-Learning Model}
The MobileNetV2 backbone pre-trained on ImageNet was loaded with
\texttt{include\_top=False} and \texttt{weights='imagenet'}, removing
the original 1000-class classification head~\cite{mobilenetv2}.
The backbone weights were initially frozen (\texttt{trainable=False})
to preserve the pre-trained representations during the early training
phase.
A \texttt{GlobalAveragePooling2D} layer was added to reduce spatial
feature maps to a single vector.
A \texttt{Dense} layer with 128 units and ReLU activation provides
domain-specific feature transformation.
A \texttt{Dropout} layer with rate 0.5 prevents overfitting.
The output layer is a single sigmoid unit for binary classification.
The full model was compiled with binary cross-entropy loss and the
Adam optimizer at a learning rate of $1 \times 10^{-4}$.
The total parameter count is 2{,}422{,}081, of which 164{,}097 are
trainable and 2{,}257{,}984 are frozen backbone weights.

\subsubsection{Feature Extraction and Fine-Tuning Strategy}
Both models used the Adam optimizer at an identical initial learning
rate to ensure a fair comparison.
For MobileNetV2, the backbone was entirely frozen in the initial
training phase, implementing pure feature extraction.
Fine-tuning with partial backbone unfreezing was not performed in the
current study and is identified as a direction for future work.

%%---------------------------------------------------------------------------
\subsection{Training Procedure}
%%---------------------------------------------------------------------------

\subsubsection{Loss Function}
Binary cross-entropy was selected as the loss function for both models,
defined as:
\begin{equation}
  \mathcal{L} = -\frac{1}{N}\sum_{i=1}^{N}
    \bigl[y_i \log(\hat{y}_i) + (1-y_i)\log(1-\hat{y}_i)\bigr]
  \label{Eq:BinaryCrossEntropy}
\end{equation}
where $y_i \in \{0,1\}$ is the true label for sample $i$,
$\hat{y}_i \in (0,1)$ is the sigmoid output of the model,
and $N$ is the total number of samples in the batch.
Binary cross-entropy is appropriate because the task involves
mutually exclusive binary labels.

\subsubsection{Optimizer and Learning Rate}
The Adam optimizer~\cite{tensorflow} was used for both models with an
initial learning rate of $\alpha = 1 \times 10^{-4}$.
Adam combines adaptive per-parameter learning rates with momentum
estimates, providing robust convergence on medical imaging tasks with
relatively small datasets.
No learning-rate schedule or weight decay was applied in the base
experiment.

\subsubsection{Training Controls}
The custom CNN was trained for 10 epochs with a batch size of 32.
MobileNetV2 was trained for 5 epochs with a batch size of 32.
No early stopping was applied; training was run to the maximum epoch
count.
Models were evaluated on the validation generator at each epoch.
A fixed random seed was used to ensure reproducibility.

\subsubsection{Hyperparameter Settings}
Table~\ref{tab:Hyperparams} summarizes the complete hyperparameter
configuration for both models.

\begin{table}[!ht]
\centering
\caption{Hyperparameter configuration for the custom CNN and
MobileNetV2 models. Both models share identical optimizer settings
to ensure a fair comparative evaluation.}
\label{tab:Hyperparams}
\begin{tabular}{|l|c|c|}
\hline
\textbf{Hyperparameter} & \textbf{Custom CNN} & \textbf{MobileNetV2} \\
\hline
Input size              & $224 \times 224 \times 3$ & $224 \times 224 \times 3$ \\
Batch size              & 32        & 32        \\
Epochs                  & 10        & 5         \\
Optimizer               & Adam      & Adam      \\
Learning rate           & $1\times10^{-4}$ & $1\times10^{-4}$ \\
Loss function           & Binary cross-entropy & Binary cross-entropy \\
Output activation       & Sigmoid   & Sigmoid   \\
Dropout rate            & 0.5       & 0.5       \\
Dense units (head)      & 128       & 128       \\
Backbone frozen         & N/A       & Yes       \\
\hline
\end{tabular}
\end{table}

%%---------------------------------------------------------------------------
\subsection{Evaluation Framework}
%%---------------------------------------------------------------------------

\subsubsection{Primary Predictive Metrics}
The following metrics are computed on the held-out test set.

\textbf{Accuracy} measures the proportion of all test images correctly
classified:
\begin{equation}
  \text{Accuracy} = \frac{TP + TN}{TP + TN + FP + FN}
  \label{Eq:Accuracy}
\end{equation}
where $TP$ denotes correctly identified PNEUMONIA cases, $TN$ denotes
correctly identified NORMAL cases, $FP$ denotes NORMAL cases predicted
as PNEUMONIA, and $FN$ denotes PNEUMONIA cases predicted as NORMAL.

\textbf{Precision} measures the fraction of PNEUMONIA predictions that
are correct:
\begin{equation}
  \text{Precision} = \frac{TP}{TP + FP}
  \label{Eq:Precision}
\end{equation}

\textbf{Recall (Sensitivity)} measures the fraction of true PNEUMONIA
cases detected:
\begin{equation}
  \text{Recall} = \frac{TP}{TP + FN}
  \label{Eq:Recall}
\end{equation}

\textbf{F1-score} is the harmonic mean of precision and recall:
\begin{equation}
  F_1 = 2 \cdot \frac{\text{Precision} \times \text{Recall}}
                     {\text{Precision} + \text{Recall}}
  \label{Eq:F1}
\end{equation}

\textbf{Specificity} measures the fraction of true NORMAL cases
correctly identified:
\begin{equation}
  \text{Specificity} = \frac{TN}{TN + FP}
  \label{Eq:Specificity}
\end{equation}

In the clinical context, recall (sensitivity) for PNEUMONIA is of
particular importance, as false negatives represent missed diagnoses
with potentially life-threatening consequences.

\subsubsection{Class-Level Evaluation}
Per-class precision, recall, and F1-score are computed from the
sklearn \texttt{classification\_report}.
Confusion matrices are generated for both models to visualize
true positive, true negative, false positive, and false negative
counts.

\subsubsection{Training Behaviour}
Training and validation accuracy and loss are plotted across all
epochs to assess convergence rate, stability, and signs of overfitting
(i.e., divergence between training and validation curves).

\subsubsection{Computational Efficiency}
The total parameter count for each model is recorded from the
\texttt{model.summary()} output.
Training time per epoch and total training time are reported where
recorded.
Test loss is used as a summary measure of generalization quality
alongside accuracy.

%%---------------------------------------------------------------------------
\subsection{Software, Hardware, and Reproducibility}
%%---------------------------------------------------------------------------
All experiments were implemented in Python 3.10 using TensorFlow 2.x~\cite{tensorflow}
and Keras~\cite{keras}.
Additional libraries include NumPy, Pandas, Matplotlib, Seaborn,
and Scikit-learn.
Training was conducted on Kaggle with GPU acceleration (NVIDIA Tesla T4).
The complete code, including notebook and saved model weights, is
available in the project repository.

Model weights were saved in HDF5 format (\texttt{.h5}) for inspection
and reuse.

%%---------------------------------------------------------------------------
\subsection{Ethical and Practical Considerations}
%%---------------------------------------------------------------------------
The dataset used in this study was collected under institutional
review board approval by the original researchers and does not contain
any personally identifiable patient information~\cite{kermany2018}.
All images are anonymized chest radiographs and their use for
non-commercial research purposes is permitted under the CC BY 4.0
license.
The models developed in this study are intended as screening-support
tools only and should not replace clinical judgment by a qualified
radiologist.
False negative predictions, in which a pneumonia-positive patient is
classified as normal, carry severe clinical consequences and must be
flagged as a critical risk in any deployment scenario.
Human oversight must be maintained in all clinical applications, and
the model should be validated on locally acquired data before
deployment in a specific clinical environment.
Potential sources of demographic or equipment bias in the training
dataset should be investigated before real-world deployment.

%%===========================================================================
\section{Results}
%%===========================================================================
Results are organized by research question.
All metrics are reported on the held-out 624-image test set unless
otherwise stated.

%%---------------------------------------------------------------------------
\subsection{Dataset and Preprocessing Results (RQ1)}
%%---------------------------------------------------------------------------
The training set contains 5{,}216 images: 1{,}341 NORMAL and 3{,}875
PNEUMONIA, representing a 2.9:1 class imbalance.
The validation set contains 16 images and the test set 624 images
(234 NORMAL, 390 PNEUMONIA).
All 5{,}863 images were successfully loaded without corruption errors.
After resizing to $224 \times 224$ pixels and normalizing to $[0,1]$,
images were passed directly to the data generators.


\begin{figure}[!ht] \centering \includegraphics[width=0.7\linewidth]{Figures/class_distribution.png} \caption{Distribution of chest X-ray images in the training dataset.} \label{fig:class_distribution} \end{figure}
\begin{figure}[!ht] \centering \includegraphics[width=0.8\linewidth]{Figures/augmentation_examples.png} \caption{Examples of augmented chest X-ray images generated during training.} \label{fig:augmentation} \end{figure}

%%---------------------------------------------------------------------------
\subsection{Training and Convergence Results}
%%---------------------------------------------------------------------------
Training curves for both models are shown in Figure~\ref{fig:training_curves}.
The custom CNN was trained for 10 epochs; training accuracy rose
steadily from approximately 87\% in epoch 1 to above 91\% by epoch 3,
while validation accuracy fluctuated between 75\% and 81\% throughout,
indicating that the small 16-image validation set introduces high
variance in epoch-level monitoring.
MobileNetV2 converged more rapidly over 5 epochs, reaching training
accuracy above 93\% by epoch 3 with validation accuracy of 87.5\%,
suggesting that the pre-trained ImageNet features provide a stronger
initialization and require fewer parameter updates to adapt to the
pneumonia classification task.

\begin{figure*}[!ht] \centering  \includegraphics[width=0.48\textwidth]{Figures/cnn_accuracy_loss.png} \hfill \includegraphics[width=0.48\textwidth]{Figures/mobilenet_accuracy_loss.png}  \caption{Training and validation accuracy and loss curves of the Custom CNN (left) and MobileNetV2 (right).} \label{fig:training_curves} \end{figure*}

%%---------------------------------------------------------------------------
\subsection{Overall Model Performance (RQ2)}
%%---------------------------------------------------------------------------
Table~\ref{tab:ModelComparison} presents the overall performance of
both models on the held-out test set.
MobileNetV2 outperforms the custom CNN on both accuracy and test loss,
achieving 88.46\% accuracy versus 71.47\%, and a substantially lower
test loss of 0.2419 versus 0.5932.
The reduction in test loss of 59.2\% indicates superior calibration
and generalization for MobileNetV2.
The 17 percentage-point accuracy gap is larger than the difference in
training-set accuracy suggests, indicating that the custom CNN
overfits the training distribution and does not generalize as well to
the held-out test set.

\begin{table}[!ht]
\centering
\caption{Overall performance comparison of the custom CNN and
MobileNetV2 on the 624-image held-out test set.}
\label{tab:ModelComparison}
\begin{tabular}{|l|c|c|}
\hline
\textbf{Model} & \textbf{Test Accuracy (\%)} & \textbf{Test Loss} \\
\hline
Custom CNN   & 71.47 & 0.5932 \\
MobileNetV2  & 88.46 & 0.2419 \\
\hline
\end{tabular}
\end{table}

%%---------------------------------------------------------------------------
\subsection{Class-Level Performance (RQ3)}
%%---------------------------------------------------------------------------
Table~\ref{tab:ClassMetrics} presents per-class precision, recall, and
F1-score for both models.
Confusion matrices are shown in Figure~\ref{fig:confusion_matrices}.

The custom CNN exhibits a severe class-prediction bias: it achieves
near-perfect recall for NORMAL (0.99) but only 0.55 recall for
PNEUMONIA, meaning it misses 45\% of true pneumonia cases.
This pattern is consistent with the model collapsing toward predicting
the NORMAL class for ambiguous inputs, despite class-weight
compensation.
The overall accuracy of 71.47\% is therefore misleading as a
standalone metric: the model fails on the clinically critical class.

MobileNetV2 produces a well-balanced result. PNEUMONIA recall of 0.96
means only 4\% of true pneumonia cases are missed, a substantial
improvement over the custom CNN.
NORMAL precision of 0.92 indicates few false alarms.
The weighted F1-score of 0.88 confirms consistently strong performance
across both classes.

\begin{table}[!ht]
\centering
\caption{Per-class precision, recall, and F1-score for the custom CNN
and MobileNetV2 on the held-out test set. NORMAL = class 0;
PNEUMONIA = class 1. Values are rounded to two decimal places.}
\label{tab:ClassMetrics}
\begin{tabular}{|l|l|c|c|c|}
\hline
\textbf{Model} & \textbf{Class} &
\textbf{Precision} & \textbf{Recall} & \textbf{F1-score} \\
\hline
Custom CNN   & NORMAL    & 0.57 & 0.99 & 0.72 \\
             & PNEUMONIA & 0.99 & 0.55 & 0.71 \\
\hline
MobileNetV2  & NORMAL    & 0.92 & 0.76 & 0.83 \\
             & PNEUMONIA & 0.87 & 0.96 & 0.91 \\
\hline
\end{tabular}
\end{table}

\begin{figure*}[!ht] \centering \includegraphics[width=0.45\textwidth]{Figures/cnn_confusion_matrix.png} \hfill \includegraphics[width=0.45\textwidth]{Figures/mobilenet_confusion_matrix.png} \caption{Confusion matrices for the Custom CNN (left) and MobileNetV2 (right) on the 624-image held-out test set.} \label{fig:confusion_matrices} \end{figure*}

%%---------------------------------------------------------------------------
\subsection{Efficiency and Deployment Results (RQ4)}
%%---------------------------------------------------------------------------
Table~\ref{tab:Efficiency} summarizes the computational characteristics
of both models.
MobileNetV2 has substantially fewer total parameters (2.42M versus
11.17M), of which only 164K are trainable in the frozen-backbone
configuration.
Despite having fewer trainable parameters, MobileNetV2 achieves higher
accuracy, confirming that the ImageNet pre-trained features provide
efficient domain-relevant representations.
The smaller trainable parameter count also reduces the risk of
overfitting, consistent with the lower test loss observed.
MobileNetV2 required only 5 training epochs to converge compared with
10 for the custom CNN, resulting in a shorter overall training
duration.
These characteristics make MobileNetV2 more attractive for deployment
in resource-constrained settings where training time and model
complexity are practical constraints.

\begin{table}[!ht]
\centering
\caption{Computational efficiency comparison of the custom CNN and
MobileNetV2.}
\label{tab:Efficiency}
\begin{tabular}{|l|c|c|}
\hline
\textbf{Metric} & \textbf{Custom CNN} & \textbf{MobileNetV2} \\
\hline
Total parameters      & 11{,}169{,}601 & 2{,}422{,}081 \\
Trainable parameters  & 11{,}169{,}345 & 164{,}097     \\
Frozen parameters     & 256            & 2{,}257{,}984  \\
Training epochs       & 10             & 5              \\
Test accuracy (\%)    & 71.47          & 88.46          \\
Test loss             & 0.5932         & 0.2419         \\
\hline
\end{tabular}
\end{table}

%%===========================================================================
\section{Discussion}
%%===========================================================================

%%---------------------------------------------------------------------------
\subsection{Interpretation of Dataset and Preprocessing Findings (RQ1)}
%%---------------------------------------------------------------------------
The 2.9:1 class imbalance in the training set is the most consequential
dataset characteristic influencing model behavior in this study.
Without compensation, a model trained on this distribution can achieve
approximately 74\% accuracy simply by predicting PNEUMONIA for every
input, since PNEUMONIA accounts for 74\% of training samples.
Class-weight compensation counteracts this by multiplying the loss
contribution of each NORMAL sample by a factor of approximately 1.93,
making errors on the minority class proportionally more costly during
gradient updates.
Despite this compensation, the custom CNN still collapsed toward
majority-class prediction at test time, suggesting that class weighting
alone is insufficient when the model architecture lacks the capacity
to form a robust decision boundary on a small and imbalanced dataset.
MobileNetV2's pre-trained features appear to mitigate this effect by
providing a more discriminative initialization that the classification
head can exploit even with limited labeled data.

The 16-image validation set is a significant inherited limitation of
the original dataset structure.
With only 8 images per class, a single misclassified image moves
validation accuracy by 6.25 percentage points, introducing high
variance in epoch-level monitoring that makes early stopping or
learning-rate scheduling unreliable.
The fluctuation between 75\% and 87.5\% in validation accuracy across
epochs for both models is almost certainly attributable to this small
sample size rather than genuine generalization differences between
epochs.
Future work using this dataset should construct a larger stratified
validation split, for example by reserving 10\% of the training set,
to obtain more stable epoch-level performance estimates.

%%---------------------------------------------------------------------------
\subsection{Comparative Performance of CNN Models (RQ2)}
%%---------------------------------------------------------------------------
The 17 percentage-point accuracy advantage of MobileNetV2 over the
custom CNN (88.46\% vs.\ 71.47\%) is substantially larger than
typically reported in transfer-learning comparisons on this dataset,
and reflects the difficulty of training a CNN from random
initialization on a relatively small and imbalanced dataset.
The more striking difference is in test loss (0.2419 vs.\ 0.5932),
indicating that MobileNetV2 produces better-calibrated probability
estimates, which is clinically important for threshold selection.
The custom CNN, trained from random initialization on approximately
5{,}000 images, achieves only 71.47\% accuracy and collapses toward
NORMAL class prediction despite class-weight compensation, suggesting
that the class imbalance and limited dataset size exceed the capacity
of the architecture to generalize reliably.
The depthwise separable convolution architecture and
ImageNet-derived edge and texture detectors of MobileNetV2 appear
to capture lung opacity patterns more efficiently~\cite{mobilenetv2}.
MobileNetV2 also required only 5 training epochs to converge compared
with 10 for the custom CNN, suggesting faster adaptation to the domain.

%%---------------------------------------------------------------------------
\subsection{Class-Level Reliability and Error Patterns (RQ3)}
%%---------------------------------------------------------------------------
The most clinically significant finding is the custom CNN's low PNEUMONIA
recall of 0.55, meaning that nearly half of all true pneumonia cases
in the test set were misclassified as NORMAL.
In a clinical screening context, such false negatives represent missed
diagnoses and could delay life-saving treatment.
By contrast, MobileNetV2 achieves a PNEUMONIA recall of 0.96, reducing
the missed-pneumonia rate to approximately 4\% of cases.
The custom CNN's near-perfect NORMAL recall (0.99) alongside poor
PNEUMONIA recall indicates the model has effectively learned to predict
NORMAL for the majority of inputs, a known failure mode when class
imbalance is not adequately compensated.
MobileNetV2's more balanced precision and recall across both classes
(NORMAL: 0.92/0.76; PNEUMONIA: 0.87/0.96) indicates a more robust
decision boundary that does not collapse toward either class.

%%---------------------------------------------------------------------------
\subsection{Efficiency and Practical Deployment (RQ4)}
%%---------------------------------------------------------------------------
MobileNetV2 offers a favorable efficiency-performance trade-off compared
with the custom CNN.
Its total parameter count of 2.42 million is less than a quarter of
the custom CNN's 11.17 million, yet it achieves substantially higher
accuracy and lower loss.
The frozen backbone configuration means that only 164{,}097 parameters
required gradient updates during training, significantly reducing
memory requirements and training time.
These properties make MobileNetV2 a strong candidate for deployment in
low-resource clinical environments, including settings with limited
GPU availability or where inference must run on edge hardware.
Future work should benchmark inference latency on mobile or embedded
devices to confirm practical deployment feasibility.

%%---------------------------------------------------------------------------
\subsection{Comparison with Existing Literature}
%%---------------------------------------------------------------------------
MobileNetV2's test accuracy of 88.46\% on the original Kaggle split is
consistent with published transfer-learning results on the same dataset.
Chouhan et al.\ reported 90.0\% accuracy with fine-tuned VGG16, though
their study did not apply class-weight compensation and reported no
per-class recall figures, making the comparison partially
opaque~\cite{chouhan2020}.
Liang and Zheng achieved 89.2\% accuracy with ResNet50 using a custom
augmentation pipeline~\cite{liang2020}; the marginal gap relative to the
present MobileNetV2 result is consistent with MobileNetV2 being a
lighter architecture with fewer parameters, and suggests that the
efficiency advantage of MobileNetV2 comes at a small but acceptable
accuracy cost.
Reshan et al.\ demonstrated that DenseNet121 with class-weight
balancing achieved higher per-class AUC than models trained without
balancing on the same dataset, corroborating the present finding that
class-imbalance treatment is necessary for reliable minority-class
recall~\cite{reshan2023}.

The custom CNN's 71.47\% test accuracy is lower than most published
results for custom architectures on this dataset, which typically
report 80--85\% accuracy.
This discrepancy is likely attributable to the absence of fine-tuned
augmentation strategy and the relatively shallow three-block
architecture, which has limited capacity to learn discriminative
features from a small imbalanced dataset.
The PNEUMONIA recall of 0.55 is particularly concerning and
underscores a finding consistent with Hashmi et al., who reported
that models trained without class compensation on this dataset
frequently exhibit low sensitivity on the NORMAL class or vice
versa despite acceptable overall accuracy~\cite{hashmi2020}.
Direct numerical comparison with all cited studies is complicated by
differences in dataset splitting strategy; the present study strictly
uses the original held-out test set, whereas some published results
are reported on reshuffled or cross-validated splits that are not
directly comparable.

%%---------------------------------------------------------------------------
\subsection{Practical and Scientific Contributions}
%%---------------------------------------------------------------------------
\textbf{Scientific Contributions.}
This study provides a systematic and reproducible comparison of a
custom CNN and a transfer-learning model on the same dataset split,
preprocessing pipeline, and evaluation protocol, addressing the
inconsistency that limits cross-study comparison in the existing
literature~\cite{hashmi2020,reshan2023}.
Class-weight compensation for training-set imbalance is reported
explicitly and its effect on minority-class recall is analyzed.
The large gap in PNEUMONIA recall between the two models (0.55 versus
0.96) highlights the clinical risk of relying on overall accuracy alone
as an evaluation metric.

\textbf{Practical Contributions.}
MobileNetV2 achieves 88.46\% test accuracy with a low test loss of
0.2419 using only 5 training epochs and fewer than 165{,}000 trainable
parameters, suggesting suitability as a fast-training screening-support
tool in settings where radiologist availability is limited~\cite{chouhan2020}.
Grad-CAM and SHAP explainability analysis are identified as high-priority
extensions that would further support clinical adoption of the model.

%%---------------------------------------------------------------------------
\subsection{Ethical, Social, and Safety Implications}
%%---------------------------------------------------------------------------
The proposed system is intended as a clinical decision-support tool
and must not be used as a standalone diagnostic device without
radiologist review.
False negative predictions (missed pneumonia) are the highest-risk
error type and should trigger mandatory clinical escalation in any
deployment protocol.
The training dataset originates from a single hospital and may not
represent the demographic or equipment diversity of all clinical
environments, introducing potential generalization bias~\cite{kermany2018}.
Deployment in a new clinical site requires local validation on
site-specific images before clinical use.
Model predictions should be communicated to clinicians alongside
uncertainty information, not as deterministic labels.

%%---------------------------------------------------------------------------
\subsection{Limitations}
%%---------------------------------------------------------------------------
The study uses a single public dataset from one institution, limiting
generalizability to different scanner types, patient demographics, and
radiograph acquisition protocols.
The original validation split contains only 16 images, which is
insufficient for reliable epoch-level model selection; future work
should use a larger stratified validation split.
MobileNetV2 was evaluated only in the frozen feature-extraction
configuration; fine-tuning deeper layers could further improve
performance and is left as future work~\cite{mobilenetv2}.
No external validation dataset was used, so the model's out-of-distribution
performance remains unknown.
Grad-CAM and SHAP explainability methods were not implemented in this
study; they are proposed as a priority direction for future work to
assess whether model attention corresponds to clinically relevant lung
regions.
Statistical significance testing between model performance differences
was not conducted; the observed accuracy gap may or may not be
statistically significant given the test-set size.

%%---------------------------------------------------------------------------
\subsection{Future Research Directions}
%%---------------------------------------------------------------------------
Future work should implement Grad-CAM and SHAP explainability analysis
to assess whether MobileNetV2 attends to clinically meaningful lung
opacity regions, and to identify any evidence of shortcut learning.
MobileNetV2 fine-tuning with partial or full backbone unfreezing should
be evaluated to assess whether further performance improvements are
achievable.
External validation on a second chest X-ray dataset, such as NIH
ChestX-ray14 or CheXpert, would establish the model's
out-of-distribution robustness.
Additional transfer-learning architectures including DenseNet121~\cite{huang2017densenet},
EfficientNetB0, and ResNet50~\cite{he2016resnet} should be compared within the same
experimental framework.
Multi-class extension to distinguish bacterial pneumonia, viral
pneumonia, COVID-19, and normal chest X-rays would significantly
increase clinical utility.
Domain-expert radiologist evaluation of explanation outputs would
provide clinically grounded validation of model trustworthiness.

%%===========================================================================
\section{Conclusion}
%%===========================================================================
This study addressed the problem of automatically classifying chest
X-ray images as NORMAL or PNEUMONIA using a comparative CNN-based
framework.
A custom three-block CNN was trained as a baseline and compared with a
MobileNetV2 transfer-learning model on the publicly available
Chest X-Ray Images (Pneumonia) dataset containing 5{,}863 images~\cite{kermany2018}.
MobileNetV2 achieved the superior overall performance with a test
accuracy of 88.46\% and a test loss of 0.2419, outperforming the
custom CNN (71.47\% accuracy, loss 0.5932) by 17 percentage points and
reducing test loss by 59\%.
The most clinically significant finding is MobileNetV2's PNEUMONIA
recall of 0.96, compared with 0.55 for the custom CNN, indicating that
transfer learning substantially reduces the rate of missed pneumonia
diagnoses.
The practical contribution of this work is a reproducible and
computationally efficient pipeline that could assist radiograph
screening in resource-constrained clinical settings.
The principal limitations are reliance on a single institutional
dataset with a very small validation split, and the absence of
Grad-CAM and SHAP explainability analysis, which are proposed as
priority directions for future work.

\bibliographystyle{elsarticle-num}
\bibliography{Bibliography}

\end{document}
